\chapter{Education and Social Change}

\section{Gandhi on Education (Basic Education / Nai Talim)}

\begin{KeyConcept}
\begin{itemize}
\item \textbf{Gandhi was against Western education} --- education based on book-learning, memorising and reproducing facts for marks.
\item Why? Western education \textbf{prepares people only for low clerical positions in the government}, and it \textbf{creates inequality}: the English-educated start looking down on the cobbler and carpenter (occupations that were respected before the British came). So Western education \textbf{destroys indigenous knowledge and talent} (carpentry, Ayurveda) and makes people feel superior.
\item His solution --- \textbf{Basic Education / Nai Talim (new training)}: introduce \textbf{productive handicrafts} (carpentry, pottery, spinning) at the school level for everyone, thereby \textbf{restoring dignity and honour to people in such occupations}. (Spinning the charkha, for instance, teaches the history of cotton while doing practical work --- bookish knowledge is useless.)
\item \textbf{Basic education is the agent of social transformation} --- train everyone in handicraft so that, when they grow up, they treat manual workers as equals.
\end{itemize}
\end{KeyConcept}

\section{Tagore on Education}

\begin{KeyConcept}
\begin{itemize}
\item \textbf{Tagore was also against Western education}: the true spirit of education lies in \textbf{self-actualization, intellectual development and critical thinking} (critical analysis, evaluative skills, logical reasoning).
\item He proposed the \textbf{mother tongue as the medium of instruction} (English should not be compulsory --- it is just a language; critical thinking does not need English).
\item He favoured teaching \textbf{history and geography through field work} --- visiting museums, forts (with their light-and-sound shows) --- rather than learning from books (e.g. 'page 200 to page 300 will be tested tomorrow').
\end{itemize}
\end{KeyConcept}

\section{Durkheim and the Functions of Education}

\begin{KeyConcept}
\begin{itemize}
\item \textbf{Durkheim} (a functionalist) looked at the \emph{functions} of education (not what education \emph{should} be):
\item 1. \textbf{Value consensus}: teaching history and ethics creates a feeling of 'us' among students from different regions (a Southerner, a Kashmiri, a Northeaster) --- cultivating \textbf{fraternity, nationalism, cooperation, teamwork and morality} --- integrating the different parts of society under one value consensus.
\item 2. \textbf{Meritocracy}: the more educated you are, the more you are westernised, and the more you look at the world through a humanitarian, egalitarian --- hence meritocratic --- outlook.
\item The \textbf{school is the agent of socialization} that transmits these functions, and their result is \textbf{social solidarity} (specialization $\rightarrow$ interdependence $\rightarrow$ solidarity).
\end{itemize}
\end{KeyConcept}

\begin{KeyConcept}
\begin{itemize}
\item \textbf{Talcott Parsons} (also functionalist): the \textbf{school is a bridge between the family and society}. The \textbf{family} socialises on \textbf{particularistic, ascriptive} lines (pattern variable A), while the \textbf{school and society} instil \textbf{universalistic, achievement-oriented, meritocratic} values (pattern variable B).
\item Education \textbf{creates a pool of talented individuals / a productive workforce}, helping in \textbf{role allocation and the division of labour} --- with the talented filtered out through \textbf{examination} and \textbf{certification}.
\end{itemize}
\end{KeyConcept}

\begin{QuickRevision}
\begin{itemize}
\item Gandhi: against Western education (book-learning, clerical jobs, inequality, kills indigenous talent); Basic Education/Nai Talim (productive handicrafts at school; dignity to manual occupations; agent of social transformation).
\item Tagore: self-actualization + critical thinking; mother tongue as medium; field-work/museums, not books.
\item Durkheim: value consensus (history/ethics $\rightarrow$ fraternity/nationalism/morality) + meritocracy; school = agent of socialization $\rightarrow$ solidarity.
\item Parsons: school as bridge (family = particularistic/ascriptive; school/society = universalistic/achievement); role allocation, division of labour, examination/certification.
\end{itemize}
\end{QuickRevision}

\section{The Marxist Critique: Bowles \& Gintis and Althusser}

\begin{CriticalPoint}
\begin{itemize}
\item \textbf{Bowles and Gintis --- the myth of meritocracy}: meritocracy does not exist. Equal opportunity is a myth because \textbf{achievement is greatly influenced by class, caste and gender position rather than individual ability}. The functionalist myth says IQ $\rightarrow$ education $\rightarrow$ class; the reality is \textbf{class $\rightarrow$ education $\rightarrow$ IQ} (your social location determines your education, which determines your 'IQ').
\item The \textbf{NFHS (National Family Health Survey) 2020} confirms this: \textbf{access to education is lower among Dalits, Muslims and tribals} --- so it is not just individual talent.
\item \textbf{Louis Althusser --- education as an Ideological State Apparatus (ISA)}: education \textbf{legitimates and reproduces inequalities}, producing an \textbf{obedient workforce by rewarding conformity over intelligence}. The myth of meritocracy hides these inequalities (you became an IAS officer, I became a clerk --- 'on merit', though merit does not exist).
\item Althusser's deeper point (echoing Amartya Sen): \textbf{poverty is not just lack of wealth but lack of capability} --- a bright transgender student or a Kashmiri student faces external constraints (discrimination, no internet, closed schools) that are not about individual ability.
\item The \textbf{New Education Policy (NEP) 2020} addresses this with \textbf{vocational occupations} and \textbf{Special Education Zones} for weaker sections and geographically excluded regions (Kashmir, the North-East).
\end{itemize}
\end{CriticalPoint}

\begin{QuickRevision}
\begin{itemize}
\item Bowles \& Gintis: myth of meritocracy; class $\rightarrow$ education $\rightarrow$ IQ (not IQ $\rightarrow$ education $\rightarrow$ class); NFHS 2020 (education lower among Dalits/Muslims/tribals).
\item Althusser: education = ISA; legitimates/reproduces inequality; rewards conformity over intelligence; myth of meritocracy hides inequality; poverty = lack of capability (Sen).
\item NEP 2020: vocational occupations + Special Education Zones (Kashmir, North-East).
\end{itemize}
\end{QuickRevision}

